\section{Introduction}
\label{sec:intro}

Floating-point arithmetic underlies nearly all scientific and numerical
software, yet binary64 does not provide enough precision for every application
\citep{Higham2002}.  When higher precision is needed, one option is a
multiple-precision library such as MPFR, which supports a user-selected working
precision \citep{Fousse2007}.  Double-word arithmetic is a more specialized
option: it represents each value with a pair of floating-point numbers and
provides roughly twice the precision of the underlying format
\citep{Dekker1971,Joldes2017}.  Its binary64 instance, called double-double
arithmetic, can be faster than a general-purpose multiple-precision arithmetic
library when roughly twice the precision of binary64 is sufficient \citep{Li2002}.
For example, double-double summation has been used to improve reproducibility
in parallel climate modeling \citep{HeDing2001}, and double-double arithmetic
has supplied additional precision for calculations in quantum electrodynamics
\citep{Aoyama2015}.

While double-word arithmetic offers performance advantages, using it safely
requires vigilance.  \citet{Li2002} called double-double arithmetic ``an
`attractive nuisance,' like an unfenced backyard swimming pool.''  The problem
is that, while double-word arithmetic provides roughly twice the precision of
the underlying format, it does not provide the usual guarantees of conventional
IEEE~754 arithmetic with twice the precision.  IEEE~754 specifies each basic
operation by correct rounding of its exact result \citep{IEEE754-2019}; under
the standard model, this gives each basic arithmetic operation the same
relative-error bound \citep{Higham2002}.  In contrast, for double-word
arithmetic, error bounds can differ across operations and even between
algorithms for the same operation.  Each algorithm must also be shown to
produce a valid double-word representation so that its output can be used as
an input to later operations.  Without this representation property,
double-word operations cannot be used soundly in sequence.  Whether an
algorithm for a double-word operation meets these two requirements depends on
conditions over its inputs and intermediate values.

To enable the safe use of double-word arithmetic, we put a verified fence
around C implementations of double-word operations.  The fence is a guarded C
interface for each double-word operation.  Before calling the arithmetic
kernel, the interface checks a sufficient, operation-specific condition.  Our
numerical analysis proves that any result produced under this condition
satisfies an operation-specific error bound and is a valid double-word
representation.  If the check fails, the kernel is not called and no numerical
result is returned.  If the function returns normally, its contract guarantees
that its result satisfies the error bound proved for that operation and is a
valid double-word representation.  The conditions required for safe
double-word arithmetic are therefore enforced by the library rather than left
to callers.

Our numerical analysis builds on existing work.  \citet{Joldes2017} established
tight error bounds for basic double-word algorithms.  \citet{Rideau2022} later
mechanized these analyses in Rocq using Flocq.  Both analyses assume an
unbounded exponent range; the mechanization represents this assumption using
Flocq's FLX format \citep{Boldo2011}.  Within FLX, these analyses establish
operation-specific error bounds and valid double-word results under their
stated hypotheses.  These guarantees do not carry directly to IEEE~754
arithmetic because FLX omits overflow and gradual underflow.  Applying them to
IEEE~754 arithmetic requires ruling out overflow and examining each rounded
intermediate that may underflow.  Where IEEE~754 rounding agrees with FLX, the
existing analysis applies.  Where it differs, the IEEE~754 behavior must be
analyzed directly.  The difference can affect the provable error bound, which
may have a larger relative-error coefficient or include an additional absolute
error term. Finally, the existing results also concern numerical algorithms rather
than their C implementations and do not provide executable checks for their
conditions.

We address the bounded-format problem independently of the C implementation.
For each operation, we define a Rocq function that specifies the sequence of
rounded floating-point operations in the numerical algorithm; we call this
function its \emph{functional model}.  We then introduce a new,
operation-specific sufficient condition that combines the hypotheses inherited
from the FLX analysis with the range requirements of a bounded-exponent format,
including conditions that rule out overflow.  Using Flocq's FLT rounding model,
which captures gradual underflow through a lower exponent bound, we prove that
whenever this condition holds, the functional model satisfies an
operation-specific error bound and produces a valid double-word representation.
These results are parameterized by the precision and exponent bounds of the
underlying radix-two format, subject to the stated format assumptions.  The
verified C library instantiates them for binary64.

We then use the Verified Software Toolchain \citep{Appel2014} to connect these
numerical results to the C implementations.  For each operation, a VST
refinement proof shows that the C kernel produces the result specified by its
binary64 functional model.
We separately prove the runtime check sound: if the check succeeds, the
operation-specific condition holds.  The checks may conservatively reject
inputs.  The proof of the guarded entry point composes these results.  If the
check succeeds, its postcondition establishes the condition and the wrapper
calls the kernel.  If the check fails, the wrapper follows its specified
nonreturning error path without calling the kernel.  For a normally returning
call, the refinement proof identifies the returned pair with the functional-model
result, and the numerical theorem establishes that this result satisfies the
operation-specific error bound and is a valid double-word representation.  The
guarded function's VST contract exposes these guarantees to C callers.

We apply this structure to \TwoSum{}, \FastTwoSum{}, and \FastTwoMult{}, which
are building blocks of double-word arithmetic, and to addition, multiplication,
and division with either one double-word and one binary64 operand or two
double-word operands.  The machine-checked results establish functional
correctness, operation-specific error bounds, and validity of the returned
double-word representations for normally returning calls at the C interface
modeled by VST.  These guarantees for both the runtime checks and the
double-word operations require compilation that preserves the strict
IEEE~754 semantics assumed by the proofs.  This requirement applies to any
compiler used to build the library and is satisfied by the CompCert compiler.
If CompCert is used, its correctness theorem additionally guarantees that the
verified C behavior is preserved in the generated assembly.

\paragraph{Contributions.}
This paper makes the following contributions:
\begin{itemize}
  \item operation-specific conditions and guarded C interfaces that enforce
        them through verified runtime checks.  The contract for a normally
        returning call guarantees the proved error bound and a valid
        double-word representation (\Cref{sec:overview,sec:specification});

  \item Rocq and Flocq proofs showing that each operation-specific condition
        implies an error bound and that the result is a valid double-word
        representation.  These results are parameterized by precision and
        exponent bounds and include conditions preventing overflow and bounds
        accounting for gradual underflow (\Cref{sec:ideas});

  \item VST proofs that the C kernels implement their functional models, the
        runtime checks establish the required conditions, and the guarded
        entry points provide the resulting contracts through the linked C
        interface (\Cref{sec:mechanization,sec:case-study}); and

  \item an evaluation of verification coverage, guarded-interface overhead,
        and the input domains accepted by the checks, with QD as a reference
        implementation (\Cref{sec:evaluation}).
\end{itemize}

The accompanying artifact contains the C sources, Rocq and VST proofs, and
evaluation scripts at \url{ARTIFACT-URL}, revision
\texttt{ARTIFACT-COMMIT}.  We begin by following one guarded multiplication
from its C interface to its numerical contract.  We then separate the
numerical and program-verification obligations, show how their results are
connected across the library, and evaluate the resulting coverage, accepted
input domains, and runtime cost.